% Part: sets-functions-relations
% Chapter: relations
% Section: orders

\documentclass[../../../include/open-logic-section]{subfiles}

\begin{document}

\olfileid{sfr}{rel}{ord}
\olsection{क्रम}

\begin{explain}
अनेक तुलनांमध्ये आपण काही वस्तू इतर वस्तूंपेक्षा एखाद्या दृष्टीने
``लहान'', ``समान'' किंवा ``मोठ्या'' आहेत असे म्हणतो. यात \emph{क्रम}
संबंध येतात. पण क्रमसंबंधांचे वेगवेगळे प्रकार आहेत. उदाहरणार्थ, काहींमध्ये
कोणत्याही दोन वस्तूंची तुलना करता आली पाहिजे, तर इतरांमध्ये तशी अट नसते.
काहींमध्ये एकरूपतेचा समावेश असतो (जसे $\le$), तर काहींमध्ये नसतो
(जसे $<$). येथे या प्रकारांचे वर्गीकरण उपयोगी पडेल.
\end{explain}

\begin{defn}[पूर्वक्रम]
परावर्ती आणि संक्रमक असणाऱ्या संबंधाला \emph{पूर्वक्रम} म्हणतात.
\end{defn}

\begin{defn}[अंशतः क्रम]
प्रतिसममितही असणाऱ्या पूर्वक्रमाला \emph{अंशतः क्रम} म्हणतात.
\end{defn}

\begin{defn}[रेषीय क्रम]\ollabel{def:linearorder}
संयोजितही असणाऱ्या अंशतः क्रमाला \emph{पूर्ण क्रम} किंवा \emph{रेषीय क्रम}
म्हणतात.
\end{defn}

\begin{ex}
प्रत्येक रेषीय क्रम अंशतः क्रमही असतो आणि प्रत्येक अंशतः क्रम पूर्वक्रमही
असतो; पण यांची उलट विधाने सत्य नाहीत. $A$ वरील सार्वत्रिक संबंध परावर्ती
आणि संक्रमक असल्याने तो पूर्वक्रम असतो. पण $A$ मध्ये एकाहून अधिक
!!{element} असतील, तर सार्वत्रिक संबंध प्रतिसममित नसतो, म्हणून तो
अंशतः क्रम नसतो.
\end{ex}

\begin{ex}
$\Bin^*$ वरील \emph{अधिक लांब नसण्याचा} $\preccurlyeq$ हा संबंध पाहा:
$x \preccurlyeq y$ तेव्हा आणि केवळ तेव्हाच, जेव्हा $\len{x} \le \len{y}$.
हा पूर्वक्रम (परावर्ती आणि संक्रमक) आहे, शिवाय संयोजितही आहे; पण
प्रतिसममित नसल्याने अंशतः क्रम नाही. उदाहरणार्थ,
$01 \preccurlyeq 10$ आणि $10 \preccurlyeq 01$, पण $01 \neq 10$.
\end{ex}

\begin{ex}
संचांच्या संचावरील $\subseteq$ हा एक महत्त्वाचा अंशतः क्रम आहे.
सर्वसाधारणपणे तो रेषीय क्रम नसतो. कारण $a \neq b$ असेल आणि
$\Pow{\{a, b\}} = \{\emptyset, \{a\}, \{b\}, \{a,b\}\}$ विचारात घेतला,
तर $\{a\} \nsubseteq \{b\}$, $\{a\} \neq \{b\}$ आणि
$\{b\} \nsubseteq \{a\}$ असे दिसते.
\end{ex}

\begin{ex}
\emph{निःशेष विभाज्यतेचा} संबंध हा रेषीय नसलेला अंशतः क्रम देतो.
$n$ आणि $m$ या पूर्णांकांसाठी, $n \mid m$ म्हणजे $n$ ने $m$ ला निःशेष
भाग जातो, अर्थात एखादा पूर्णांक $k$ असा असतो की $m = kn$.
$\Nat$ वर हा अंशतः क्रम आहे, पण रेषीय क्रम नाही: उदाहरणार्थ,
$2 \nmid 3$ आणि $3 \nmid 2$. $\Int$ वरील संबंध म्हणून पाहिल्यास
विभाज्यता केवळ पूर्वक्रम आहे, कारण ती प्रतिसममित नाही:
$1 \mid -1$ आणि $-1 \mid 1$, पण $1 \neq -1$.
\end{ex}

\begin{ex}
अनुक्रमांच्या $A^*$ या संचावरील \emph{विस्तार} संबंध असा आहे:
$s \sqsubseteq s'$ तेव्हा आणि केवळ तेव्हाच, जेव्हा $s = \emptyseq$
(रिक्त अनुक्रम), $s = s'$, किंवा $s = \tuple{s_1, \dots, s_n}$ आणि
$s' = \tuple{s_1, \dots, s_n, s_{n+1}, \dots, s_m}$.
$s \sqsubseteq s'$ असल्यास $s$ हा $s'$ चा \emph{आरंभीचा खंड} आहे असेही
म्हणतो. $A^*$ वरील विस्तार संबंध अंशतः क्रम आहे, पण रेषीय क्रम नाही;
उदा., $a \neq b$ असेल, तर $ab \not\sqsubseteq ba$ आणि
$ba \not\sqsubseteq ab$.
\end{ex}

\begin{defn}[काटेकोर क्रम]
\emph{काटेकोर क्रम} म्हणजे अपरावर्ती, असममित आणि संक्रमक असणारा संबंध.
\end{defn}

\begin{defn}[काटेकोर रेषीय क्रम]\ollabel{def:strictlinearorder}
संयोजितही असणाऱ्या काटेकोर क्रमाला \emph{काटेकोर पूर्ण क्रम} किंवा
\emph{काटेकोर रेषीय क्रम} म्हणतात.
\end{defn}

\begin{ex}
$\le$ हा काटेकोर रेषीय क्रम $<$ याला अनुरूप रेषीय क्रम आहे.
$\subseteq$ हा काटेकोर क्रम $\subsetneq$ याला अनुरूप अंशतः क्रम आहे.
\end{ex}

$A$ वरील कोणत्याही काटेकोर क्रम $R$ मध्ये कर्ण $\Id{A}$, म्हणजे सर्व
$\tuple{x, x}$ जोड्या, मिळवल्यास अंशतः क्रम तयार होतो. (याला $R$ चे
\emph{परावर्ती संवरण} म्हणतात.) उलट, अंशतः क्रमातून $\Id{A}$ काढून
टाकल्यास काटेकोर क्रम मिळतो. पुढील दोन निष्कर्ष हे नेमके स्पष्ट करतात.

\begin{prop}\ollabel{prop:stricttopartial}
$R$ हा $A$ वरील काटेकोर क्रम असेल, तर $R^+ = R \cup \Id{A}$ हा अंशतः
क्रम असतो. शिवाय, $R$ हा काटेकोर रेषीय क्रम असेल, तर $R^+$ हा रेषीय
क्रम असतो.
\end{prop}

\begin{proof}
$R$ हा काटेकोर क्रम आहे असे समजा, म्हणजे $R \subseteq A^2$ आणि $R$
अपरावर्ती, असममित व संक्रमक आहे. $R^+ = R \cup \Id{A}$ असू द्या.
$R^+$ परावर्ती, प्रतिसममित आणि संक्रमक आहे हे दाखवायचे आहे.

$R^+$ स्पष्टपणे परावर्ती आहे, कारण सर्व $x \in A$ साठी
$\tuple{x, x} \in \Id{A} \subseteq R^+$.

$R^+$ प्रतिसममित आहे हे दाखवण्यासाठी व्याघात काढू. $R^+xy$ आणि $R^+yx$
असूनही $x \neq y$ आहे असे समजा. $\tuple{x,y} \in R \cup \Id{A}$,
पण $\tuple{x, y} \notin \Id{A}$ असल्याने $\tuple{x, y} \in R$, म्हणजे
$Rxy$ असले पाहिजे. त्याचप्रमाणे $Ryx$. पण $R$ असममित आहे या
गृहीतकाशी हे विसंगत आहे.

संक्रमकता दाखवण्यासाठी $R^+xy$ आणि $R^+yz$ असे समजा.
$\tuple{x, y} \in R$ आणि $\tuple{y,z} \in R$ दोन्ही असतील, तर $R$
संक्रमक असल्याने $\tuple{x, z} \in R$. अन्यथा $\tuple{x, y} \in \Id{A}$,
म्हणजे $x = y$, किंवा $\tuple{y, z} \in \Id{A}$, म्हणजे $y = z$.
पहिल्या प्रकरणात गृहीतकानुसार $R^+yz$ आणि $x = y$, म्हणून $R^+xz$.
दुसऱ्या प्रकरणातही त्याचप्रमाणे. प्रत्येक प्रकरणात $R^+xz$, म्हणून
$R^+$ संक्रमकही आहे.

``शिवाय'' या भागासाठी, $R$ संयोजितही आहे असे समजा. म्हणून सर्व
$x \neq y$ साठी $Rxy$ किंवा $Ryx$, म्हणजे $\tuple{x, y} \in R$ किंवा
$\tuple{y, x} \in R$. $R \subseteq R^+$ असल्याने $R^+$ बाबतही हे सत्य
राहते, म्हणून $R^+$ संयोजितही आहे.
\end{proof}

\begin{prop}\ollabel{prop:partialtostrict}
$R$ हा $A$ वरील अंशतः क्रम असेल, तर $R^- = R \setminus \Id{A}$ हा
काटेकोर क्रम असतो. शिवाय, $R$ हा रेषीय क्रम असेल, तर $R^-$ हा काटेकोर
रेषीय क्रम असतो.
\end{prop}

\begin{proof}
हे सरावासाठी ठेवले आहे.
\end{proof}

\begin{prob}
\olref[sfr][rel][ord]{prop:partialtostrict} याची सिद्धता द्या.
\end{prob}

काटेकोर रेषीय क्रम विस्तारात्मकतेसारखा एक गुणधर्म पाळतात, हे पुढील
सोप्या निष्कर्षातून दिसते:

\begin{prop}\ollabel{prop:extensionality-strictlinearorders}
$<$ हा $A$ वरील काटेकोर रेषीय क्रम असेल, तर:
\[
  (\forall a, b \in A)((\forall x \in A)(x < a \liff x < b) \lif a = b).
\]
\end{prop}

\begin{proof}
$(\forall x \in A)(x < a \liff x < b)$ असे समजा. $a < b$ असेल, तर
$a < a$; पण $<$ अपरावर्ती आहे याच्याशी हे विसंगत आहे. म्हणून $a \nless b$.
अगदी त्याचप्रमाणे $b \nless a$. $<$ संयोजित असल्याने $a = b$.
\end{proof}

\end{document}
